% 来源：OEBPS/Text/01_47.xhtml（OPF spine 顺序；仅第一册）
\cleardoublepage
\BookEntry[理查德·费曼《别闹了，费曼先生！》]{理查德·费曼《别闹了，费曼先生！》}

\BookCover{images/书单31别闹了，费曼先生！.jpg}{[美] 理查德·费曼著| 吴程远译| 生活·读书·新知三联书店2005年版}

\begin{BookSummaryBox}
\textbf{诺贝尔物理学奖得主费曼的“顽童趣事”}
\par
\end{BookSummaryBox}

\subsubsection*{这本书谈什么？}

理查德·费曼先生是诺贝尔物理学奖得主，他近代理论物理方面最有造诣的科学家，却又是一个幽默机智、几近顽童的人。这本书便讲述了他生活中的种种乐趣。
\par

他是加州理工学院物理系的教授，20世纪30年代在普林斯顿大学毕业后，被征召加入制造原子弹的曼哈顿计划。但他生性好奇，且颇为机智，在森严的安保系统监控下，却以破解安全锁为乐，取得机密资料后，留下便签告诫政府要小心。他爱坐在酒吧做科学研究，也曾在酒吧被取缔后亲自到法庭上辩护。他几乎不参与加州理工学院内各项行政工作，为了自由，他连续五年辞去美国国家科学院院士的荣誉位置，因为选举其他院士的责任颇困扰他。而他也可能是历史上唯一被按摩院请去画裸体画、在酒吧厕所内与醉鬼大打出手的科学家。他曾经跟爱因斯坦和波耳等大师讨论物理问题，也曾在赌城中与职业赌徒研究输赢的机率！
\par

\subsubsection*{芒格为什么谈起这本书？}

在芒格演讲“专业人士需要更多的跨学科技能”中，他提出：
\par

\begin{quotation}
软科学应该加强模仿硬科学的基本治学精神与方法。（这里所说的硬科学，是指数学、物理学、化学和工程学四门学科。）这种治学精神与方法值得效仿。毕竟，硬科学在如下两方面做得更好：（1）避免单一学科造成的错误；（2）使得大量跨学科知识更容易被应用，并时常获得良好的结果，比如物理学家理查德·费曼就能够用一只橡皮环解释挑战者号航天飞机爆炸的原因。而且以前软科学也曾借鉴这种治学精神与方法，取得了很好的效果。
\par
\end{quotation}

理查德·费曼出生在纽约的法洛卡威。他本科毕业于麻省理工学院物理学系，从普林斯顿大学得到博士学位。他参与了曼哈顿计划，对原子弹的研发起到了重要作用。他一直在康奈尔大学任教，然后在1951年跳槽到加州理工学院。费曼对物理学的主要贡献是量子电动力学、电磁放射、原子和其他更为基本的粒子之间的关系。1965年，他和施温格、朝永振一郎共同获得诺贝尔物理学奖。
\par

晚年的费曼受委托对“挑战者号”航天飞机事故进行调查。他展示了橡皮环如何遇冷萎缩，导致高温的燃气外泄，从而引起了爆炸。
\par

\subsubsection*{《别闹了，费曼先生！》作者理查德·费曼如是说}

“冯·诺依曼教会了我一个很有趣的想法：你不需要为身处的世界负任何责任。因此我就形成了强烈的‘社会不负责任感’，从此成为一个快活逍遥的人。大家听好了，我的不负责任感全都是由于冯·诺依曼在我思想上撒下的种子而起的。”
\par

“在课堂上时，你可以思考一些已经很清楚的基本东西。这些知识有趣、令人愉快，重温一遍又何妨？另一方面，有没有更好的介绍方式？有什么相关的新问题？你能不能赋予这些旧知识新生命？基本的东西思考起来并不难；而如果你没想出什么新东西来，没关系，以前想过的已足以应付讲课之用了。但如果你真的有什么新想法，能从新角度看事物，你会觉得很愉快。”
\par

“学生问的问题，有时也能提供新的研究方向。他们经常提出一些我曾经思考过、但暂时放弃、却都是些意义很深远的问题。重新想想这些问题，看看能否有所突破，也很有意思。学生未必理解我想回答的方向，或者是我想思考的层次；但他们问我这个问题，却往往提醒了我相关的问题。单单靠自己，是不容易获得这种启示的。”
\par

“因此对我来说，教学与学生们使我的生命继续发光发亮，我也永远不会接受任何人替我安排一切。”
\par
